    \begin{tikzpicture}[
        neuron/.style={
            circle, 
            draw=black!60, 
            very thick, 
            fill=white, 
            inner sep=0pt,
            minimum size=4pt
        },
        output/.style={
            circle, 
            draw=black, 
            very thick, 
            fill=white, 
            inner sep=0pt,
            minimum size=20pt
        },
        concept/.style={
            circle,
            draw=black,
            very thick,
            inner sep=0pt,
            minimum size=20pt
        },
        c_text/.style={
            rectangle, 
            rounded corners, 
            draw=black, 
            fill=pearTwo!30,
            thick, 
            anchor=west
        },
        >=stealth
    ]

    \draw[fill=red!50, rounded corners, dashed, very thick, opacity=0.6] 
        (-2.5, 2) rectangle (3, -2);
    \draw[fill=yellow!50, rounded corners, dashed, very thick, opacity=0.6] 
        (4, 2) rectangle (9, -2);
    \draw[fill=blue!50,rounded corners, dashed, very thick, opacity=0.6] 
        (1.7, 2.2) rectangle (7.1, -2.2);
    
    \node[fill=red!50, rectangle, anchor=center] 
        (exp1) at (-0.3, 2.7) {\small CRL on unlabelled data};
    \node[fill=blue!50, rectangle, anchor=center] 
        (exp2) at (4.5, 2.7) {\small Learning alignments with concept labels};
    \node[fill=yellow!50, rectangle, anchor=center] 
        (exp3) at (8.7, 2.7) {\small Label predictor};

    \node[
        rectangle, 
        rounded corners, 
        draw=black, 
        very thick, 
        fill=purple!70]
        (pic) at (-1.1, 0) 
        {\includegraphics[scale=0.8]{./figs/causal_ident/teapot.png}};

    % Input layer (3 neurons)
    \foreach \i in {1,2,3} {
        \node[neuron] (I\i) at (0,0.6-0.3*\i) {};
        \draw[->] (pic) -- (I\i);
    }

    % Hidden layer (6 neurons)
    \foreach \i in {1,2,3,4,5,6} {
        \node[neuron] (H\i) at (0.5,1.05-0.3*\i) {};
    }

    % Hidden layer (6 neurons)
    \foreach \i in {1,2,3,4,5,6} {
        \node[neuron] (H_2\i) at (1,1.05-0.3*\i) {};
    }

    % Output layer (4 neurons)
    \foreach \i in {1,2,3,4} {
        \node[output, fill=blue!50] (O\i) at (2.5,2.5 - 1*\i) {$M_\i$};
    }

    \foreach \i in {1,2,3,4} {
        \node[concept, fill=red!50] (C\i) at (4.5,2.5-1*\i) {$C_\i$};
    }

    \node[c_text]
        (concept_1) at (5.2, 1.5) {\scriptsize Color object};
    \node[c_text]
        (concept_2) at (5.2, 0.5) {\scriptsize Rot.\ object};   
    \node[c_text]
        (concept_3) at (5.2,-0.5) {\scriptsize Rot.\ spotlight};
    \node[c_text]
        (concept_4) at (5.2,-1.5) {\scriptsize Object};

    \node[draw=black, circle, very thick, fill=orange] 
        (Y) at (8, 0) {\large $Y$};

    % Connections between input and hidden layer
    \foreach \i in {1,2,3} {
        \foreach \j in {1,2,3,4,5,6} {
            \draw[->, opacity=0.4] (I\i) -- (H\j);
        }
    }

    % Connections between input and hidden layer
    \foreach \i in {1,2,3,4,5,6} {
        \foreach \j in {1,2,3,4,5,6} {
            \draw[->, opacity=0.4] (H\i) -- (H_2\j);
        }
    }

    % Connections between hidden and output layer
    \foreach \i in {1,2,3,4,5,6} {
        \foreach \j in {1,2,3,4} {
            \draw[->, opacity=0.4] (H_2\i) -- (O\j);
        }
    }

    \foreach \i in {1, 2, 3, 4} {
        \draw[->, thick] (concept_\i) -- (Y);
    }  

    \draw[->, dashed, color=red, very thick] (O1) -- (C3);
    \draw[->, dashed, color=red, very thick] (O2) -- (C4);
    \draw[->, dashed, color=red, very thick] (O3) -- (C2);
    \draw[->, dashed, color=red, very thick] (O4) -- (C1);


    \node (alpha) at ($(O4)!0.5!(C4)$) {\large $\alpha$};
    \draw [shorten >= 0.2cm] (alpha) -- ($(O2)!0.55!(C4)$);
\end{tikzpicture}
